\documentclass[10pt,a4paper]{amsart}
\usepackage{amsmath,amssymb,amsthm,graphicx}
\usepackage{verbatim}
\newenvironment{reference}{\comment}{\endcomment}
%\newenvironment{reference}{}{}
\newenvironment{slogan}{\comment}{\endcomment}
\newenvironment{history}{\comment}{\endcomment}

% For commutative diagrams we use Xy-pic
\usepackage[all]{xy}

% We use 2cell for 2-commutative diagrams.
\xyoption{2cell}
% Theorem environments.
%
\theoremstyle{plain}
\newtheorem{theorem}[subsection]{Theorem}
\newtheorem{proposition}[subsection]{Proposition}
\newtheorem{lemma}[subsection]{Lemma}

\theoremstyle{definition}
\newtheorem{definition}[subsection]{Definition}
\newtheorem{example}[subsection]{Example}
\newtheorem{exercise}[subsection]{Exercise}
\newtheorem{situation}[subsection]{Situation}

\theoremstyle{remark}
\newtheorem{remark}[subsection]{Remark}
\newtheorem{remarks}[subsection]{Remarks}

\numberwithin{equation}{subsection}

% Macros
%
\def\lim{\mathop{\mathrm{lim}}\nolimits}
\def\colim{\mathop{\mathrm{colim}}\nolimits}
\def\Spec{\mathop{\mathrm{Spec}}}
\def\Hom{\mathop{\mathrm{Hom}}\nolimits}
\def\Ext{\mathop{\mathrm{Ext}}\nolimits}
\def\SheafHom{\mathop{\mathcal{H}\!\mathit{om}}\nolimits}
\def\SheafExt{\mathop{\mathcal{E}\!\mathit{xt}}\nolimits}
\def\Sch{\mathit{Sch}}
\def\Mor{\mathop{\mathrm{Mor}}\nolimits}
\def\Ob{\mathop{\mathrm{Ob}}\nolimits}
\def\Sh{\mathop{\mathit{Sh}}\nolimits}
\def\NL{\mathop{N\!L}\nolimits}
\def\CH{\mathop{\mathrm{CH}}\nolimits}
\def\proetale{{pro\text{-}\acute{e}tale}}
\def\etale{{\acute{e}tale}}
\def\QCoh{\mathit{QCoh}}
\def\Ker{\mathop{\mathrm{Ker}}}
\def\Im{\mathop{\mathrm{Im}}}
\def\Coker{\mathop{\mathrm{Coker}}}
\def\Coim{\mathop{\mathrm{Coim}}}

% Boxtimes
%
\DeclareMathSymbol{\boxtimes}{\mathbin}{AMSa}{"02}

%
% Macros for moduli stacks/spaces
%
\def\QCohstack{\mathcal{QC}\!\mathit{oh}}
\def\Cohstack{\mathcal{C}\!\mathit{oh}}
\def\Spacesstack{\mathcal{S}\!\mathit{paces}}
\def\Quotfunctor{\mathrm{Quot}}
\def\Hilbfunctor{\mathrm{Hilb}}
\def\Curvesstack{\mathcal{C}\!\mathit{urves}}
\def\Polarizedstack{\mathcal{P}\!\mathit{olarized}}
\def\Complexesstack{\mathcal{C}\!\mathit{omplexes}}
% \Pic is the operator that assigns to X its picard group, usage \Pic(X)
% \Picardstack_{X/B} denotes the Picard stack of X over B
% \Picardfunctor_{X/B} denotes the Picard functor of X over B
\def\Pic{\mathop{\mathrm{Pic}}\nolimits}
\def\Picardstack{\mathcal{P}\!\mathit{ic}}
\def\Picardfunctor{\mathrm{Pic}}
\def\Deformationcategory{\mathcal{D}\!\mathit{ef}}
\usepackage{cleveref}
\crefname{originaldefinition}{original definition}{original definitions}
\Crefname{originaldefinition}{Original definition}{Original definitions}
\makeatletter
\newlabel{definition-valuative-criterion}{{03IX}{source}{}{sourcepointer-1}{}}
\newlabel{definition-valuative-criterion@cref}{{[originaldefinition][0][]03IX}{[1][0][]source}{}{}{}}
\makeatother
\crefname{originallemma}{original lemma}{original lemmas}
\Crefname{originallemma}{Original lemma}{Original lemmas}
\makeatletter
\newlabel{descent-lemma-descending-property-etale}{{02VN}{source}{}{sourcepointer-2}{}}
\newlabel{descent-lemma-descending-property-etale@cref}{{[originallemma][0][]02VN}{[1][0][]source}{}{}{}}
\makeatother
\crefname{originallemma}{original lemma}{original lemmas}
\Crefname{originallemma}{Original lemma}{Original lemmas}
\makeatletter
\newlabel{lemma-base-change-separated}{{03KL}{source}{}{sourcepointer-3}{}}
\newlabel{lemma-base-change-separated@cref}{{[originallemma][0][]03KL}{[1][0][]source}{}{}{}}
\makeatother
\crefname{originallemma}{original lemma}{original lemmas}
\Crefname{originallemma}{Original lemma}{Original lemmas}
\makeatletter
\newlabel{lemma-flat-base-change-scheme-theoretic-image}{{089E}{source}{}{sourcepointer-4}{}}
\newlabel{lemma-flat-base-change-scheme-theoretic-image@cref}{{[originallemma][0][]089E}{[1][0][]source}{}{}{}}
\makeatother
\crefname{originallemma}{original lemma}{original lemmas}
\Crefname{originallemma}{Original lemma}{Original lemmas}
\makeatletter
\newlabel{lemma-quasi-compact-permanence}{{03KS}{source}{}{sourcepointer-5}{}}
\newlabel{lemma-quasi-compact-permanence@cref}{{[originallemma][0][]03KS}{[1][0][]source}{}{}{}}
\makeatother
\crefname{originallemma}{original lemma}{original lemmas}
\Crefname{originallemma}{Original lemma}{Original lemmas}
\makeatletter
\newlabel{lemma-quasi-compact-scheme-theoretic-image}{{082Z}{source}{}{sourcepointer-6}{}}
\newlabel{lemma-quasi-compact-scheme-theoretic-image@cref}{{[originallemma][0][]082Z}{[1][0][]source}{}{}{}}
\makeatother
\crefname{originallemma}{original lemma}{original lemmas}
\Crefname{originallemma}{Original lemma}{Original lemmas}
\makeatletter
\newlabel{lemma-scheme-theoretic-image-reduced}{{0830}{source}{}{sourcepointer-7}{}}
\newlabel{lemma-scheme-theoretic-image-reduced@cref}{{[originallemma][0][]0830}{[1][0][]source}{}{}{}}
\makeatother
\crefname{originallemma}{original lemma}{original lemmas}
\Crefname{originallemma}{Original lemma}{Original lemmas}
\makeatletter
\newlabel{lemma-section-immersion}{{03KP}{source}{}{sourcepointer-8}{}}
\newlabel{lemma-section-immersion@cref}{{[originallemma][0][]03KP}{[1][0][]source}{}{}{}}
\makeatother
\crefname{originallemma}{original lemma}{original lemmas}
\Crefname{originallemma}{Original lemma}{Original lemmas}
\makeatletter
\newlabel{more-algebra-lemma-valuation-ring-torsion-free-flat}{{0539}{source}{}{sourcepointer-9}{}}
\newlabel{more-algebra-lemma-valuation-ring-torsion-free-flat@cref}{{[originallemma][0][]0539}{[1][0][]source}{}{}{}}
\makeatother
\crefname{originallemma}{original lemma}{original lemmas}
\Crefname{originallemma}{Original lemma}{Original lemmas}
\makeatletter
\newlabel{spaces-lemma-space-presentation}{{0262}{source}{}{sourcepointer-10}{}}
\newlabel{spaces-lemma-space-presentation@cref}{{[originallemma][0][]0262}{[1][0][]source}{}{}{}}
\makeatother
\crefname{originallemma}{original lemma}{original lemmas}
\Crefname{originallemma}{Original lemma}{Original lemmas}
\makeatletter
\newlabel{spaces-properties-lemma-quasi-compact-affine-cover}{{03H6}{source}{}{sourcepointer-11}{}}
\newlabel{spaces-properties-lemma-quasi-compact-affine-cover@cref}{{[originallemma][0][]03H6}{[1][0][]source}{}{}{}}
\makeatother
\title{Pushing down a solution of the valuative criterion along a separated morphism of algebraic spaces}
\author{The Stacks Project Authors}
\date{stacks-project:spaces-morphisms.tex:local-snapshot-20261012; GNU FDL 1.2 or later}
\begin{document}
\maketitle
\setcounter{subsection}{1}
\noindent\textit{Editorial note.} All delivered mathematical text is the author's own, copied from the pinned Stacks Project chapter file spaces-morphisms.tex; no statement or proof step was rewritten, reordered or removed.\par\medskip
\noindent\textit{Editorial note.} The counted statement refers to the diagram of Definition definition-valuative-criterion; that definition is a navigation pointer here and the argument only uses the commutativity of the displayed diagram, which is part of the statement span.\par\medskip
\noindent\textit{Editorial note.} The eleven references of the counted spans are recorded as non-load-bearing navigation pointers; the argument that consumes them stays inside the counted proof span.\par\medskip
\noindent\textit{Editorial note.} Preamble: the three native excerpt spans of sources/stacks-spaces-morphisms/preamble.tex (lines 8-13, 14-18 and 154-217). The chapter file itself carries no preamble text.\par\medskip
\noindent\textit{Allowed standard prerequisites (editorial).}\begin{itemize}
\item Base change of a separated morphism of algebraic spaces and the shape of the valuative diagram. Lemma lemma-base-change-separated and Definition definition-valuative-criterion, used at proof lines 8530-8549. Reduces to the case where the base is Spec(A) with A a valuation ring.
\item Scheme theoretic image of a quasi-compact morphism and its base change by a flat morphism. Lemmas lemma-scheme-theoretic-image-reduced and lemma-flat-base-change-scheme-theoretic-image, used at proof lines 8556-8570. Identifies Z\_K with the base change of the scheme theoretic image Z.
\item A section of a separated morphism is a closed immersion. Lemma lemma-section-immersion, used at proof lines 8573-8576. Turns the given Spec(K) -> Z\_K into a closed immersion and yields Z\_K = Spec(K).
\item Torsion free modules over a valuation ring are flat. More on Algebra, Lemma more-algebra-lemma-valuation-ring-torsion-free-flat, used at proof lines 8588-8593. Converts torsion freeness of B and B tensor\_A B over A into flatness.
\item Descent of etaleness along a faithfully flat ring map. Descent, Lemma descent-lemma-descending-property-etale, used at proof lines 8603-8607. Concludes that A -> B is etale, hence that the quotient Z is Spec(A).
\end{itemize}\medskip
\ifdefined\hypertarget\hypertarget{sourcepointer-1}{}\fi
\noindent\textit{Source reference.} Original definition 03IX: focused-seeds/AGP/sources/stacks-snapshot-20261011/tags.txt, source lines 4619--4619. Navigation pointer into the Stacks Project; the standard fact it supplies is applied inside the counted proof span and the argument resting on it stays there.\par\medskip
\ifdefined\hypertarget\hypertarget{sourcepointer-2}{}\fi
\noindent\textit{Source reference.} Original lemma 02VN: focused-seeds/AGP/sources/stacks-snapshot-20261011/tags.txt, source lines 3774--3774. Navigation pointer into the Stacks Project; the standard fact it supplies is applied inside the counted proof span and the argument resting on it stays there.\par\medskip
\ifdefined\hypertarget\hypertarget{sourcepointer-3}{}\fi
\noindent\textit{Source reference.} Original lemma 03KL: focused-seeds/AGP/sources/stacks-snapshot-20261011/tags.txt, source lines 4681--4681. Navigation pointer into the Stacks Project; the standard fact it supplies is applied inside the counted proof span and the argument resting on it stays there.\par\medskip
\ifdefined\hypertarget\hypertarget{sourcepointer-4}{}\fi
\noindent\textit{Source reference.} Original lemma 089E: focused-seeds/AGP/sources/stacks-snapshot-20261011/tags.txt, source lines 10494--10494. Navigation pointer into the Stacks Project; the standard fact it supplies is applied inside the counted proof span and the argument resting on it stays there.\par\medskip
\ifdefined\hypertarget\hypertarget{sourcepointer-5}{}\fi
\noindent\textit{Source reference.} Original lemma 03KS: focused-seeds/AGP/sources/stacks-snapshot-20261011/tags.txt, source lines 4688--4688. Navigation pointer into the Stacks Project; the standard fact it supplies is applied inside the counted proof span and the argument resting on it stays there.\par\medskip
\ifdefined\hypertarget\hypertarget{sourcepointer-6}{}\fi
\noindent\textit{Source reference.} Original lemma 082Z: focused-seeds/AGP/sources/stacks-snapshot-20261011/tags.txt, source lines 10269--10269. Navigation pointer into the Stacks Project; the standard fact it supplies is applied inside the counted proof span and the argument resting on it stays there.\par\medskip
\ifdefined\hypertarget\hypertarget{sourcepointer-7}{}\fi
\noindent\textit{Source reference.} Original lemma 0830: focused-seeds/AGP/sources/stacks-snapshot-20261011/tags.txt, source lines 10270--10270. Navigation pointer into the Stacks Project; the standard fact it supplies is applied inside the counted proof span and the argument resting on it stays there.\par\medskip
\ifdefined\hypertarget\hypertarget{sourcepointer-8}{}\fi
\noindent\textit{Source reference.} Original lemma 03KP: focused-seeds/AGP/sources/stacks-snapshot-20261011/tags.txt, source lines 4685--4685. Navigation pointer into the Stacks Project; the standard fact it supplies is applied inside the counted proof span and the argument resting on it stays there.\par\medskip
\ifdefined\hypertarget\hypertarget{sourcepointer-9}{}\fi
\noindent\textit{Source reference.} Original lemma 0539: focused-seeds/AGP/sources/stacks-snapshot-20261011/tags.txt, source lines 6597--6597. Navigation pointer into the Stacks Project; the standard fact it supplies is applied inside the counted proof span and the argument resting on it stays there.\par\medskip
\ifdefined\hypertarget\hypertarget{sourcepointer-10}{}\fi
\noindent\textit{Source reference.} Original lemma 0262: focused-seeds/AGP/sources/stacks-snapshot-20261011/tags.txt, source lines 2835--2835. Navigation pointer into the Stacks Project; the standard fact it supplies is applied inside the counted proof span and the argument resting on it stays there.\par\medskip
\ifdefined\hypertarget\hypertarget{sourcepointer-11}{}\fi
\noindent\textit{Source reference.} Original lemma 03H6: focused-seeds/AGP/sources/stacks-snapshot-20261011/tags.txt, source lines 4555--4555. Navigation pointer into the Stacks Project; the standard fact it supplies is applied inside the counted proof span and the argument resting on it stays there.\par\medskip
\begin{lemma}
\label{lemma-push-down-solution}
Let $S$ be a scheme. Let $f : X \to Y$ be a separated morphism of
algebraic spaces over $S$. Suppose given a diagram
$$
\xymatrix{
\Spec(K') \ar[r] \ar[d] & \Spec(K) \ar[r] & X \ar[d] \\
\Spec(A') \ar[r] \ar[rru] & \Spec(A) \ar[r] \ar@{-->}[ru] & Y
}
$$
as in Definition \ref{definition-valuative-criterion} with $K \subset K'$
arbitrary. Then the dotted arrow exists making the diagram commute.
\end{lemma}

\begin{proof}
We have to show that we can find a morphism $\Spec(A) \to X$ fitting
into the diagram.

\medskip\noindent
Consider the base change $X_A = \Spec(A) \times_Y X$ of $X$.
Then $X_A \to \Spec(A)$ is a separated morphism of algebraic
spaces (Lemma \ref{lemma-base-change-separated}). Base changing
all the morphisms of the diagram above we obtain
$$
\xymatrix{
\Spec(K') \ar[r] \ar[d] & \Spec(K) \ar[r] & X_A \ar[d] \\
\Spec(A') \ar[r] \ar[rru] & \Spec(A) \ar@{=}[r] & \Spec(A)
}
$$
Thus we may replace $X$ by $X_A$, assume that $Y = \Spec(A)$ and
that we have a diagram as above. We may and do replace $X$ by
a quasi-compact open subspace containing the image of $|\Spec(A')| \to |X|$.

\medskip\noindent
The morphism $\Spec(A') \to X$ is quasi-compact by
Lemma \ref{lemma-quasi-compact-permanence}.
Let $Z \subset X$ be the scheme theoretic image of
$\Spec(A') \to X$. Then $Z$ is a reduced
(Lemma \ref{lemma-scheme-theoretic-image-reduced}),
quasi-compact (as a closed subspace of $X$), separated
(as a closed subspace of $X$) algebraic space over $A$.
Consider the base change
$$
\Spec(K') = \Spec(A') \times_{\Spec(A)} \Spec(K) \to
X \times_{\Spec(A)} \Spec(K) = X_K
$$
of the morphism $\Spec(A') \to X$ by the flat morphism of schemes
$\Spec(K) \to \Spec(A)$. By
Lemma \ref{lemma-flat-base-change-scheme-theoretic-image}
we see that the scheme theoretic image of this morphism
is the base change $Z_K$ of $Z$. On the other hand, by
assumption (i.e., the commutative diagram above)
this morphism factors through a morphism
$\Spec(K) \to Z_K$ which is a section to the structure morphism
$Z_K \to \Spec(K)$. As $Z_K$ is separated, this section is a
closed immersion (Lemma \ref{lemma-section-immersion}).
We conclude that $Z_K = \Spec(K)$.

\medskip\noindent
Let $V \to Z$ be a surjective \'etale morphism with $V$ an
affine scheme (Properties of Spaces, Lemma
\ref{spaces-properties-lemma-quasi-compact-affine-cover}).
Say $V = \Spec(B)$. Then $V \times_Z \Spec(A') = \Spec(C)$
is affine as $Z$ is separated. Note that $B \to C$ is injective
as $V$ is the scheme theoretic image of $V \times_Z \Spec(A') \to V$
by Lemma \ref{lemma-quasi-compact-scheme-theoretic-image}.
On the other hand, $A' \to C$ is \'etale as corresponds
to the base change of $V \to Z$.
Since $A'$ is a torsion free $A$-module,
the flatness of $A' \to C$ implies $C$ is a torsion free
$A$-module, hence $B$ is a torsion free $A$-module.
Note that being torsion free as an $A$-module is
equivalent to being flat (More on Algebra, Lemma
\ref{more-algebra-lemma-valuation-ring-torsion-free-flat}).
Next, we write
$$
V \times_Z V = \Spec(B')
$$
Note that the two ring maps $B \to B'$ are \'etale as $V \to Z$
is \'etale. The canonical surjective map $B \otimes_A B \to B'$
becomes an isomorphism after tensoring with $K$ over $A$
because $Z_K = \Spec(K)$. However, $B \otimes_A B$ is torsion
free as an $A$-module by our remarks above. Thus $B' = B \otimes_A B$.
It follows that the base change of the ring map
$A \to B$ by the faithfully flat ring map $A \to B$
is \'etale (note that $\Spec(B) \to \Spec(A)$ is surjective
as $X \to \Spec(A)$ is surjective). Hence $A \to B$ is
\'etale (Descent, Lemma \ref{descent-lemma-descending-property-etale}),
in other words, $V \to X$ is \'etale.
Since we have $V \times_Z V = V \times_{\Spec(A)} V$
we conclude that $Z = \Spec(A)$ as algebraic spaces
(for example by Spaces, Lemma \ref{spaces-lemma-space-presentation})
and the proof is complete.
\end{proof}

\end{document}
